% Generated draft: requires mathematical and source review.
\documentclass[11pt]{article}
\usepackage[margin=25mm]{geometry}
\usepackage{fontspec}
\setmainfont{Latin Modern Roman}
\usepackage{amsmath,amssymb,mathtools,unicode-math}
\setmathfont{Latin Modern Math}
\usepackage{enumitem,needspace,xurl,hyperref}
\setlength{\parindent}{0pt}
\setlength{\parskip}{5pt}
\setlength{\emergencystretch}{3em}
\newcommand{\N}{\mathbb{N}}
\newcommand{\Z}{\mathbb{Z}}
\newcommand{\Q}{\mathbb{Q}}
\newcommand{\R}{\mathbb{R}}
\newcommand{\C}{\mathbb{C}}
\newcommand{\F}{\mathbb{F}}
\newcommand{\A}{\mathbb{A}}
\newcommand{\PP}{\mathbb{P}}
\newcommand{\E}{\mathbb{E}}
\newcommand{\eps}{\varepsilon}
\newcommand{\dd}{\,\mathrm{d}}
\begin{document}
\section*{UnsolvedMath ID: 2302014}
\textbf{Source problem code:} catalogue entry
\subsection{Definitions and mathematical statement}
For a nonzero entire function $f(z)=\sum_{n=0}^\infty a_nz^n$ and $r>0$, define its maximum term and maximum modulus by
\[
m(r)=\max_{n\geq0}|a_n|r^n,\qquad
M(r,f)=\max_{|z|=r}|f(z)|.
\]
The zero function is excluded because the quotient would be undefined. The ratio is at most $1$ by the coefficient estimate.

The supplied Problem 2.14 contains three questions:
\begin{enumerate}
\item Given $C>1/2$, does there exist an entire function $f$ with $m(r)/M(r,f)\to C$ as $r\to\infty$? The source says each $0<C\leq1/2$ is attainable. Since the ratio is at most $1$, any nontrivial unresolved parameter range must lie in $1/2<C\leq1$; the source does not explicitly state this upper restriction.
\item For a zero-free entire function, the source states $\liminf_{r\to\infty}m(r)/M(r,f)=0$ and asks whether the full limit is zero. Here ``$f(z)\ne0$'' is read as nonvanishing at every point, rather than merely not identically zero. Nonconstantness is missing: nonzero constant functions have ratio $1$ and contradict the stated preliminary liminf assertion.
\item Put $\beta_f=\liminf_{r\to\infty}m(r)/M(r,f)$. Determine the exact supremum $\beta$ over the intended class of entire functions. The source reports
\[
\frac47<\beta<\frac2\pi.
\]
\end{enumerate}
The common function class is a material unresolved qualification. As transcribed, arbitrary entire functions include nonzero polynomials, for which $m(r)/M(r,f)\to1$. Thus the displayed bounds for $\beta$ cannot describe the supremum over all nonzero entire functions. A restriction to transcendental entire functions is a natural candidate for the omitted context, but is not verified from the supplied record. If that is confirmed, the precise third target would be
\[
\beta=\sup_{f\text{ transcendental entire}}\beta_f,
\]
and the second target would concern nonconstant zero-free entire functions. These are explicitly conditional repairs, not silent changes to the source. The reported attainable values, liminf assertion, and strict bounds are retained with this domain issue visible. The supplied summary reports the exact values unresolved as of the 2018 edition, with no recent resolution located.
\subsection{Short English statement}
The source studies limiting ratios of the largest Taylor-series term to the maximum modulus, asks about nonvanishing functions, and asks for an exact extremal liminf. Its reported bounds require an omitted restriction excluding polynomials; the zero-free subquestion also needs nonconstantness. Those domain issues are preserved rather than silently repaired.
\subsection{Sources}
\begin{itemize}
\item Local UnsolvedMath registry, permanent ID 2302014, exact source code catalogue entry, ``Research Problems in Function Theory --- Problem 2.14''; status open, classification OPEN-TRIAGE.
\item Hayman, \emph{Research Problems in Function Theory} (2018), Problem 2.14, \url{https://arxiv.org/abs/1809.07200}. The record identifies extraction from public arXiv source TeX, accessed July 29, 2026; the 2018 update was retained for status triage.
\item Clunie and Hayman, cited in the source for this collection of questions and bounds, without a title, date, or complete bibliographic reference.
\item The supplied research summary says the problem was open in Hayman's 2018 edition, exact values were unknown, and no recent resolution was located. No additional citation or precise search date is supplied.
\end{itemize}
\end{document}
